% !PRINTED 189
\weberpage{189}

From this $a_{m,n}$ can be found, once the earlier $a_{m,n}$, that is those
in which $2m+3n$ is smaller than $\nu$, have already been calculated.
% [sic] German source prints an ungrammatical relative construction; English conveys the intended relation.

One sees from this that the $a_{m,n}$ contain only powers of 3 in the
denominator. In the \glqq Formeln und Lehrsätzen zum Gebrauch der
elliptischen Funktionen\grqq{} edited by H. A. \gesperrt{Schwarz}, after
lectures of \gesperrt{Weierstrass} (2. Ausgabe 1893), the $a_{m,n}$ have
been calculated up to $\nu=17$. It appears thereby that these coefficients
are not only integers, but that they also contain as factors powers of 3
which increase with $\nu$. I am unable to give a proof of this.

\begin{center}
\rule{2cm}{0.4pt}
\end{center}
